\documentclass[../../main.tex]{subfiles}% 注意这里的文件路径不能用 ./main.tex ,否则用latexmk编译子文件会报错
\graphicspath{{\subfix{./image/}}{\subfix{../../image/}}} % 指定图片引用路径,后续可以直接使用图片文件名
% 如果图片存放在上级目录,则使用 ../../image/ 作为备用路径

% 例如:
% \begin{figure}[H]
% \centering
% \includegraphics[width=0.3\textwidth]{图.png}
% \caption{}
% \label{figure:图}
% \end{figure}
% 注意:上述\label{}一定要放在\caption{}之后,否则引用图片序号会只会显示??.

\begin{document}

\section{不可测集}

\begin{theorem}\label{theorem:不可测集}
设 \(\bm{Q}^n\) 为 \(\bm{R}^n\) 中的有理点集. 定义等价关系 \(x \sim y\) 当且仅当 \(x-y\in \bm{Q}^n\). 根据这一等价关系 \(\sim\),将 \(\bm{R}^n\) 中一切点分类,凡有等价关系者均属一类 (例如 \(\bm{Q}^n\) 本身即为其中一类). 根据选择公理,在每一类中取一点且只取一点形成点集 \(W\),则 \(W\) 为不可测集.
\end{theorem}
\begin{remark}
关于不可测集 \(W\),需注意上述不可测集 \(W\) 是从 \(\bm{R}^n\) 中的一切等价类中选出来形成的,但实际上并未明确界定如何从每类中选出一个元素,也就是说选择方法有很多,因此 \(W\) 可以有无穷多种. 若把任一种选择法而形成的 \(W\) 记成选成集,并记一切选成集全体为 \(S\),则对一切 \(\alpha > 0\),均存在 \(W\in S\),使得 \(m^*(W)=\alpha\).
\end{remark}
\begin{proof}
假设 \(W\) 为可测集,则第一种情形是 \(m(W) > 0\). 由\cref{theorem:定理2.19}和\hyperref[theorem:Steinhaus定理]{Steinhaus定理},可知点集 \(W \setminus W\) 含有一球 \(B(0, \delta)\),因此存在 \(x\in (W \setminus W)\cap \bm{Q}^n\) 且 \(x\neq 0\). 即存在 \(y,z\in W\) 使得 \(x=y-z\),与 \(W\) 的构造矛盾. 第二种情形 \(m(W)=0\) 也不能发生,若作可列个平移集 \(W+\{r^{(k)}\}\),其中 \(\{ r^{(1)}, r^{(2)}, \cdots, r^{(k)}, \cdots \} = \bm{Q}^n\). 显然有
\begin{align}
\bm{R}^n = \bigcup\limits_{k=1}^\infty (W+\{r^{(k)}\}). \label{eq:1}
\end{align}
再根据 \eqref{eq:1},因为 \(m(W)=0\),从而 \(m(W+\{r^{(k)}\})=0\),进而 \(\bm{R}^n\) 是零测集,矛盾!所以 \(W\) 不是Lebesgue可测集.
\end{proof}

比较两个点集的内含大或多,还是小或少,其度量方法有多种: 用基数是一种,用测度也是一种,还有一种观点是用所谓“第一纲集”和“第二纲集”,若用后者来看可测集,则“大多数”具有正外测度的点集都是不可测的.

若 \(W\) 是不可测集,则其不良性质可从下述命题获悉: 存在 \(\varepsilon_0 > 0\),使得对满足 \(A \supseteq W\) 以及 \(B \supseteq W^c\) 的任意两个可测集 \(A\) 与 \(B\),均有 \(m(A\cap B) \geqslant \varepsilon_0\). 这说明 \(W\) 与 \(W^c\) 中的点相互无休止地缠绕在一起,无法厘清. 从另一角度看,若记 \(B^c = C\),则 \(C \subseteq W\). 上述结论说明,不可测集是无法用可测集从“外包”和“内挤”两侧来逼近的.

\begin{proposition}
\begin{enumerate}[(1)]
\item \(\bm{R}^n\) 中存在零测集 \(E\),使 \(E+E\) 是不可测集.
\item \([0,1]\) 中存在不可数集 \(E\),使 \(E \setminus E\) 无内点.
\item \((0,1]\) 中存在互不相交的不可测集列 \(\{W_n\}\),使得 \(\bigcup\limits_{n=1}^\infty W_n = (0,1]\).
\item 设 \((0, +\infty) = E_1 \cup E_2\),且 \(E_1 \cap E_2 = \varnothing\). 若每个 \(E_i (i=1,2)\) (其元素)对加法运算是封闭的,则 \(E_1, E_2\) 均为不可测集.
\end{enumerate}
\end{proposition}

\begin{example}
\begin{enumerate}[(1)]
\item 证明:不存在可测集 \(E \subseteq [0,1]\),使得对于任意的 \(x\in \bm{R}\),存在 \(y\in E\),使得 \(x-y\in \bm{Q}\).
\end{enumerate}
\end{example}
\begin{proof}
假设存在这样的可测集 \(E\)。由题意可知 \(E + \bm{Q} = \bm{R}\)，即
\begin{align}
\bm{R} = \bigcup\limits_{q \in \bm{Q}} (E+q). \label{eq:unioon}
\end{align}
若 \(m(E)=0\)，则由 \eqref{eq:unioon} 可得 \(\bm{R}\) 是零测集，矛盾，因此 \(m(E) > 0\)。
根据\hyperref[theorem:Steinhaus定理]{Steinhaus定理}，必存在 \(\delta > 0\)，使得 \((-\delta, \delta) \subseteq E \setminus E\)。
取 \(q_0 \in \bm{Q} \cap (0, \delta)\)，则 \(q_0 \in E \setminus E\)。于是存在 \(y_1, y_2 \in E\)，使得 \(y_1 - y_2 = q_0\)。
由于 \(q_0 \neq 0\)，故 \(y_1 \neq y_2\)，且 \(y_1\) 与 \(y_2\) 属于同一个 \(\bm{Q}\) 等价类。
然而，要使得 \(E + \bm{Q} = \bm{R}\) 成立，\(E\) 必须在 \(\bm{R}\) 的每一个等价类中取且仅取一个代表元，此时 \(E\) 即为 \(\bm{R}\) 中的一个 Vitali集。
众所周知，Vitali集是Lebesgue不可测的，这与 \(E\) 是可测集的假设相矛盾。
因此，不存在这样的可测集 \(E \subseteq [0,1]\)。
\end{proof}







































\end{document}